\documentclass{article}
\usepackage[UTF8]{ctex}
\usepackage{amsmath,amssymb,amsfonts}
\usepackage{enumitem}
\usepackage{geometry}
\geometry{a4paper, margin=1in}
% === 错题本排版优化（不参与题目标识） ===
\setlength{\parindent}{0pt}
\setlength{\parskip}{0.35em plus 0.10em minus 0.08em}
\setlength{\emergencystretch}{3em}
\setlist{itemsep=0.25em, topsep=0.40em, parsep=0pt}
\allowdisplaybreaks[2]
% === 错题本排版优化结束 ===
\begin{document}
\section*{哈尔滨工程大学2025年数学分析真题}
\section*{85．哈尔滨工程大学2025年数学分析试卷（回忆版）}
\begin{itemize}

\subsection{题目1}
  \item[1．] 计算 \(\lim _{x \rightarrow 0}\left(\frac{1}{x^{2}}-\frac{1}{(\tan x)^{2}}\right)\) ．

\subsection{题目2}
  \item[2．] 证明：对任意的 \(x>0\) ，有 \(0<\frac{1}{\ln (1+x)}-\frac{1}{x}<1\) ．
\end{itemize}

\subsection{题目3}
3．讨论下列级数的敛散性 \(\sum_{n=3}^{\infty} \frac{1}{n(\ln n)^{p}(\ln \ln n)^{q}}\) ，其中 \(p, q\) 为实数。

\subsection{题目4}
4．设 \(f(x)\) 在 \(x=0\) 处连续，且对任意的 \(x, y \in \mathbb{R}\) ，有 \(f(x+y)=f(x)+f(y)\) ．证明：

\begin{itemize}
  \item[（1）] \(f\) 在 \(\mathbb{R}\) 上连续．
  \item[（2）] 存在实数 \(M\) ，使得 \(|f(x)| \leq M|x|\) ．
\end{itemize}

\subsection{题目5}
5．求积分 \(\int_{0}^{1} \frac{\ln x}{\sqrt{1-x^{2}}} \mathrm{~d} x\) ．

\subsection{题目6}
6．设函数 \(u(x, y)=\sum_{n=1}^{\infty} \sum_{m=1}^{\infty} A_{m n} \sin (m \pi x) \sin (n \pi y)\) 满足 \(u_{x x}+u_{y y}=x,(x, y) \in[0,1] \times[0,1]\) ，求 \(A_{m n}\) ．

\begin{itemize}

\subsection{题目7}
\end{itemize}

\subsection{题目7}
原题边界“7．”：令函数 \(f\) 为定义在 \(\mathbb{R}\) 上的连续函数，证明：对任意的 \(x, y, z \in \mathbb{R}^{+}\)，下式成立：
\[
\int_{0}^{x} \int_{0}^{y} \int_{0}^{z} f(t) \mathrm{d} t \mathrm{~d} z \mathrm{~d} y=\frac{1}{2} \int_{0}^{x}(x-t)^{2} f(t) \mathrm{d} t .
\]

\subsection{题目8}

\subsection{题目8}
原题边界“8．”：求函数 \(f(x, y)=\left(x^{2}+2 y^{2}\right) e^{-\left(x^{2}+y^{2}\right)}\) 在区域 \(D=\left\{(x, y) \mid x^{2}+y^{2} \leq 4\right\}\) 上的最大值和最小值．

\subsection{题目9}
9．计算二重积分 \(\iint_{R}\left(x^{2}+y^{2}\right) \mathrm{d} x \mathrm{~d} y\) ，其中 \(R=\left\{(x, y) \mid x^{2}-x y+y^{2} \leq 2\right\}\) ．

\subsection{题目10}
10．设 \(B=B(0,1)\) 是 \(\mathbb{R}^{3}\) 中以原点为中心以 1 为半径的开球体，\(\partial B\) 为其表面．令

\[
u=(u(x, y, z), v(x, y, z), w(x, y, z))
\]

为 \(\mathbb{R}^{3}\) 中的向量场，满足 \(u_{x}+v_{y}+w_{z}=0\) 。且在 \(\partial B\) 中满足 \(\mathbf{u} \cdot \mathbf{n}=0\) ，其中 \(\mathbf{n}\) 为 \(\partial B\) 的单位外法向量．证明下式成立：

\[
\iiint_{B}[(\mathbf{u} \cdot \nabla) \mathbf{u}] \cdot \mathbf{u} \mathrm{d} x \mathrm{~d} y \mathrm{~d} z=0 .
\]

其中 \(\mathbf{u} \cdot \nabla=u \frac{\partial}{\partial x}+v \frac{\partial}{\partial y}+w \frac{\partial}{\partial z}\) ．
\end{document}
